\documentclass[12pt]{article}

% --- Paquetes Esenciales ---
\usepackage[utf8]{inputenc}    % Caracteres en español
\usepackage[spanish]{babel}   % Idioma español
\usepackage{amsmath}          % Matemáticas avanzadas
\usepackage{geometry}         % Márgenes
\geometry{letterpaper, margin=1in}

% --- Paquete para Dibujar Circuitos Eléctricos ---
\usepackage[siunitx, american]{circuitikz} 

\begin{document}

\section{EJEMPLO DE CIRCUITOS ELÉCTRICOS}

\subsection*{CIRCUITO 1: FUENTE DC Y RESISTENCIAS}
\begin{center}
\begin{circuitikz}[scale=1.2]
    \draw (0,0) to[battery2, l=$5\text{V}$] (0,3) % Fuente DC
          to[R, l=$220\,\Omega$] (3,3)          % Resistencia
          -- (3,0) -- (0,0);                    % Cierre de malla
\end{circuitikz}
\end{center}

\subsection*{CIRCUITO 2: RESISTENCIA Y CAPACITOR EN SERIE}
\begin{center}
\begin{circuitikz}[scale=1.2]
    \draw (0,0) to[sinusoidal voltage source, l=$v(t)$] (0,3) % Fuente AC
          to[R, l=$220\,\Omega$] (2.5,3)                     % Resistencia
          to[C, l=$4\text{F}$] (5,3)                         % Capacitor en serie
          -- (5,0) -- (0,0);                                 % Cierre de malla
\end{circuitikz}
\end{center}

\subsection{CIRCUITO 3: INDUCTOR Y RESISTENCIA EN SERIE}
\begin{center}
\begin{circuitikz}[scale=1.2]
    \draw (0,0) to[battery2, l=$5\text{V}$] (0,3) % Fuente DC
          to[L, l=$5\text{mH}$] (2.5,3)           % Inductor
          to[R, l=$220\,\Omega$] (5,3)            % Resistencia en serie
          -- (5,0) -- (0,0);                      % Cierre de malla
\end{circuitikz}
\end{center}

\subsection{CIRCUITO 4: LEDS EN PARALELO}
\begin{center}
\begin{circuitikz}[scale=1.2]
    \draw (0,0) to[battery2, l=$5\text{V}$] (0,3) % Fuente DC
          -- (2,3) to[full led, l=LED1] (2,0)     % Primer LED en paralelo
          -- (0,0);
    \draw (2,3) -- (4,3) to[full led, l=LED2] (4,0) % Segundo LED en paralelo
          -- (2,0);
\end{circuitikz}
\end{center}

\end{document}